% Source: Wald GR, physical PDF p. 201
%         (printed p. 188).
% Coverage: Chapter 8 opening material before Section 8.1.
% Figures/tables/footnotes: no figures, tables, or footnotes.
% Status: source-reviewed and compile-clean.
% Uncertainties: none; a printed-source typo is corrected below.

% Wald numbers results within each subsection (e.g., Theorem 8.1.2) and restarts
% footnotes at the opening of the chapter.
\begingroup
\edef\chaptereightpreviousfootnote{\number\value{footnote}}
\setcounter{footnote}{0}
\renewcommand{\thetheorem}{\thesubsection.\arabic{theorem}}

The causal structure of spacetime in special relativity already was described briefly
in section 1.2. Associated with each event, \(p\), in spacetime is a light cone, as
illustrated in Fig.~\ref{fig:1.2}. We assign the label ``future'' to half of the cone and the label
``past'' to the other half. The events lying in the interior of the future light cone
represent events which can be reached by a material particle starting at \(p\); these
comprise the ``chronological future'' of \(p\). The chronological future of \(p\) together with
the events lying on the cone itself comprise the ``causal future'' of \(p\), which physically
represents events which, in principle, can be influenced by a signal emitted from \(p\).

In general relativity, the causal structure of spacetime is \emph{locally} of the same
qualitative nature as in the flat spacetime of special relativity. However, significant
differences can occur globally because of nontrivial topology, spacetime singularities,
or the ``twisting'' of the directions of light cones as one moves from point to point. The
purpose of this chapter is to give an account of the definitions and basic results
concerning the causal structure of spacetimes in general relativity. These results not
only are of interest in their own right but also are a crucial ingredient in the proof of
the singularity theorems, which we shall discuss in the next chapter. Further discussion
of topics on causal structure can be found in Hawking and Ellis (1973), Penrose (1972),
and Geroch (1970b); most of the discussion presented here is based on these references.

The definitions of chronological and causal futures in general spacetimes are given
in section 8.1 and several properties of these sets are derived. The conditions that
express the notion that a spacetime be ``causally well behaved'' are discussed in
section 8.2. Finally, in section 8.3 the notion of domains of dependence and global
hyperbolicity are defined and numerous properties of globally hyperbolic spacetimes
are derived. All the arguments given in this chapter rely heavily on the machinery
of topological spaces outlined in appendix A, and it is assumed that the reader has
read (or will refer to) this appendix.

The discussion throughout this chapter will concern arbitrary spacetimes \((M,g_{ab})\)
in the sense that we shall not attempt to impose Einstein's equation on \(g_{ab}\). Some of
the examples given here may appear to be artificial in that they are constructed by
removing points and/or making topological identifications of Minkowski spacetime.
Nevertheless, they provide excellent illustrations of the types of phenomena that can
occur in much less artificial spacetime models in general relativity.
% SOURCE ERRATUM: printed ``artifical.''
